\subsubsection{Receive Pending Documents}
Receive Pending Documents use case enables the registrar’s office to verify and modify the document statuses that were previously pending due to being either missing or incomplete on the part of the applicants. After the user logs in and navigates to their dashboard, the pending documents module can be accessed. It will allow the registrar to find an applicant and mark their documents as having been received. This way, all documents can be updated and verified accordingly. This use case is often utilized in the later admission process stages.
\begin{table}[H]
    \centering
    \small
    \renewcommand{\arraystretch}{1.4}
    \caption{Use Case Table: Receive Pending Documents (UC-5.11)} % ক্যাপশন উপরে আনা হয়েছে
    \label{tab:receive-pending-docs-registrar}
    \begin{tabularx}{\textwidth}{|l|X|}
        \hline
        \rowcolor[gray]{0.9} {Characteristics} & {Description} \\ \hline
        Use Case ID & UC-5.11 \\ \hline
        Use Case Name & Receive Pending Documents \\ \hline
        Primary Actors & Registrar Office \\ \hline
        Goal & Confirm and update the status of previously missing or pending documents in the system. \\ \hline
        Precondition & User must be signed in with Registrar Office credentials. \\ \hline
        Scenario & 1. Log in to the dashboard. 2. Navigate to the 'Pending Documents' section via the sidebar. 3. Locate the applicant and click the 'OK' or 'Receive' action button. \\ \hline
        Exception & Unauthorized access or document verification error. \\ \hline
        Frequency of Use & High frequency during the finalization of the admission process. \\ \hline
    \end{tabularx}
\end{table}
\FloatBarrier